\chapter*{مقدمة}
\addcontentsline{toc}{chapter}{مقدمة}
\markboth{مقدمة}{}

\begin{tcolorbox}[enhanced, colback=brandDark, colframe=brandDark, arc=4pt,
    coltext=white, left=4mm, right=4mm]
\faUserTie\quad \textbf{إعداد:} \teacher \hfill \faUsers\quad \textbf{الفئة:} \audience
\end{tcolorbox}

يجمع هذا الكتاب \textbf{\en{650} تمرينًا محلولًا} في أساسيات الخوارزميات ولغة \en{C}،
موزعة على الفصول الثلاثة عشر للمقياس. كل التمارين مستوحاة من مجالين قريبين من
الواقع المهني: \textbf{التجارة الإلكترونية} (سلات، طلبيات، مخزون، فواتير، توصيل) و\textbf{علم البيانات}
(متوسطات، تباين، تطبيع، انحدار، تصنيف، تجميع).

\section*{بنية كل تمرين}

\begin{center}
\begin{tikzpicture}[node distance=4mm,
    blk/.style={draw=#1, fill=#1!12, very thick, rounded corners=5pt,
        minimum height=1.1cm, text width=2.3cm, align=center,
        font=\khatt\bfseries\small}]
    \node[blk=brandDark] (a) {\faFile*\\ نص التمرين};
    \node[blk=brandTeal, left=of a] (b) {\faListOl\\ خطوات الحل};
    \node[blk=tealDeep, left=of b] (c) {\faCode\\ برنامج \en{C}};
    \node[blk=brandMuted, left=of c] (d) {\faTerminal\\ نتيجة التنفيذ};
    \node[blk=brandOrange, left=of d] (e) {\faProjectDiagram\\ المخطط الانسيابي};
    \foreach \x/\y in {a/b, b/c, c/d, d/e} \draw[sahm] (\x) -- (\y);
\end{tikzpicture}
\end{center}

\begin{itemize}
    \item \textbf{نص التمرين}: المسألة كما يطرحها متجر أو محلل بيانات.
    \item \textbf{خطوات الحل}: الفكرة الخوارزمية مرتبة خطوة بخطوة قبل أي كود.
    \item \textbf{البرنامج}: برنامج \en{C} كامل، تعليقاته \textcolor{cmtRed}{\textbf{باللون الأحمر}} تشرح كل سطر مهم.
    \item \textbf{نتيجة التنفيذ}: المخرجات \textbf{الحقيقية} للبرنامج (كل البرامج تُرجمت ونُفّذت).
    \item \textbf{المخطط الانسيابي} \en{(Organigramme)}: التمثيل البياني للخوارزمية.
\end{itemize}

\section*{رموز المخطط الانسيابي}

\begin{center}
\begin{tikzpicture}
    \node[fterm] (t) at (0,0) {بداية / نهاية};
    \node[anchor=east, text width=8cm] at (-2.2,0) {بداية الخوارزمية أو نهايتها.};
    \node[fio, text width=2.2cm] (i) at (0,-1.3) {\textbf{اقرأ} \fc{x}};
    \node[anchor=east, text width=8cm] at (-2.2,-1.3) {عملية إدخال (قراءة) أو إخراج (كتابة).};
    \node[fproc, text width=2.2cm] (p) at (0,-2.6) {\fc{s = s + x}};
    \node[anchor=east, text width=8cm] at (-2.2,-2.6) {معالجة: إسناد أو حساب.};
    \node[fcall, text width=2.2cm] (c) at (0,-3.9) {\fc{f(x)}};
    \node[anchor=east, text width=8cm] at (-2.2,-3.9) {استدعاء دالة (إجراء معرَّف في مكان آخر).};
    \node[fdec, text width=1.4cm] (d) at (0,-5.4) {\fc{x > 0}};
    \node[anchor=east, text width=8cm] at (-2.2,-5.4) {اختبار شرط: مخرجان \textbf{نعم} و\textbf{لا}.};
    \draw[sahm] (1.6,-6.6) -- (-0.4,-6.6);
    \node[anchor=east, text width=8cm] at (-2.2,-6.6) {اتجاه سير التنفيذ.};
\end{tikzpicture}
\end{center}

\section*{قاعدة الكتاب: لا دوال جاهزة}

\begin{obox}{\faBan\quad ممنوع استعمال الدوال المعرفة مسبقًا}
الهدف هو تعلم الخوارزميات، لذلك \textbf{نكتب بأنفسنا} كل ما توفره المكتبات عادة:
طول سلسلة، مقارنة سلسلتين، الجذر التربيعي، القيمة المطلقة، الترتيب، البحث، التحويل بين الأعداد والنصوص...
\end{obox}

\begin{center}
\begin{tabularx}{0.95\linewidth}{|>{\centering\arraybackslash}p{3.6cm}|X|}
\hline
\ras{الصنف} & \ras{الدوال} \\ \hline
\textcolor{lvEasy}{\faCheck}\ مسموح & الإدخال والإخراج: \code{printf}، \code{scanf}، \code{getchar}، \code{putchar} \\ \hline
\textcolor{lvEasy}{\faCheck}\ مسموح (الفصول \en{10--13}) & الذاكرة والملفات: \code{malloc}، \code{free}، \code{fopen}، \code{fclose}، \code{fprintf}، \code{fscanf}، \code{fgetc}، \code{fputc} \\ \hline
\textcolor{lvHard}{\faTimes}\ ممنوع & \code{strlen}، \code{strcpy}، \code{strcmp}، \code{sqrt}، \code{pow}، \code{abs}، \code{qsort}، \code{atoi}، \code{toupper}، \code{rand}، ... \\ \hline
\end{tabularx}
\end{center}

\section*{مستويات الصعوبة}

\begin{center}
\begin{tikzpicture}
\begin{axis}[ybar, bar width=1.1cm, width=11cm, height=5.5cm,
    symbolic x coords={E,M,H}, xtick=data, ymin=0, ymax=36,
    xticklabels={\textarabic{سهل}, \textarabic{متوسط}, \textarabic{صعب}},
    nodes near coords, every node near coord/.append style={font=\bfseries},
    axis lines*=left, ylabel={\textarabic{عدد التمارين في كل فصل}},
    enlarge x limits=0.3]
    \addplot[fill=brandTeal, draw=tealDeep] coordinates {(E,10) (M,30) (H,10)};
\end{axis}
\end{tikzpicture}
\end{center}

\begin{itemize}
    \item \lvbadge{1}\ تطبيق مباشر لمفهوم واحد من الفصل.
    \item \lvbadge{2}\ مسألة حقيقية تجمع عدة مفاهيم وتتطلب تفكيرًا خوارزميًا.
    \item \lvbadge{3}\ خوارزمية متقدمة (ترتيب، بحث، تجميع، تحسين...) تُبنى بالكامل يدويًا.
\end{itemize}

\section*{ترجمة البرامج وتنفيذها}

كل برنامج موجود في المجلد \code{code/chNN/exMM.c} مع نتيجة تنفيذه \code{exMM.out}:

\begin{codebox}\begin{english}
\begin{lstlisting}[style=outstyle, numbers=none]
gcc -std=c99 -Wall -Wextra code/ch04/ex12.c -o ex12
./ex12 < code/ch04/ex12.in
\end{lstlisting}
\end{english}\end{codebox}
